\documentclass[10pt]{article}

\usepackage[utf8]{inputenc}

\usepackage[T1]{fontenc}

\usepackage{amsmath}

\usepackage{amsfonts}

\usepackage{amssymb}

\usepackage[version=4]{mhchem}

\usepackage{stmaryrd}

\usepackage{bbold}

\usepackage{graphicx}

\usepackage[export]{adjustbox}

\graphicspath{ {./images/} }

\usepackage{mathrsfs}



\begin{document}

\begin{enumerate}

  \item Локальная теория. Пусть $M=\mathbb{R}^{n}$. Л а гр а нже в ы м м ногообразием в гамильтоновом (симплектическом) пространстве $T^{*} \mathbb{R}^{n}=\mathbb{R}^{n} \oplus \mathbb{R}^{n}$ с координатами

\end{enumerate}



и формой



\[

(x, p)=\left(x^{1}, \ldots, x^{n} ; p_{1}, \ldots, p_{n}\right)

\]



\[

d x \wedge d p=d x^{1} \wedge d p_{1}+\ldots+d x^{n} \wedge d p_{n}

\]



называется такое подмногообразие $L, \operatorname{dim} L=n$, для к-рого $\left.d x_{\wedge} d p\right|_{L}=0$. На лагранжевом многообразии всегда существует канонич. атлас, состоящий из карт $U_{I}$ с координатами



\[

\left(x^{I}, p_{\bar{I}}\right)=\left(x^{i_{1}}, \ldots, x^{i_{m}}, p_{i_{m+1}}, \ldots, p_{i_{n}}\right)

\]



здесь



$I=\left\{i_{1}, \ldots, i_{m}\right\}, \bar{I}=\left\{i_{m+1}, \ldots, i_{n}\right\}, I \cup \bar{I}=\{1,2, \ldots, n\}$.



Пусть $\varphi \in C_{0}^{\infty}\left(U_{I}\right)$. Локальный М. к. о. есть отображение



\[

k_{I}: C_{0}^{\infty}\left(U_{1}\right) \rightarrow C^{\infty}\left(\mathbb{R}_{x}^{n} \times \mathbb{R}_{+, h}^{1}\right),

\]



определяемое формулой



\[

k_{I} \varphi=F_{p_{\bar{I}} \rightarrow x^{1 / h}}^{\bar{I}}\left\{e^{\frac{i}{h} S_{I}\left(x^{I}, p_{\bar{I}}\right)} \sqrt{\mu_{I}\left(x^{I}, p_{\bar{I}}^{-}\right)} \varphi\left(x^{I}, p_{I}^{-}\right)\right\}

\]



Здесь $F_{p_{\bar{I}}}^{1 / h} \rightarrow x^{\bar{I}}$ есть $1 / h$-преобразование Фурье по перемен$\operatorname{HЫM} p_{\bar{I}}$ :



\[

F_{p_{\bar{I}} \rightarrow x^{\bar{I}}}^{1 / h}(f)=\left(\frac{i}{2 \pi h}\right)^{|\bar{I}| / 2} \int_{\mathbb{R}_{\bar{I}}} e^{\frac{i}{h} x^{\bar{I}} p_{\bar{I}}} f\left(p_{\bar{I}}\right) d p_{\bar{I}},

\]



$|\vec{I}|=n-m, \quad S_{I}-$ действие в карте $U_{I}$, то есть нек-рое решение уравнения $d S_{I}=p_{I} d x^{I}-\left.x^{\bar{I}} d p_{\bar{I}}\right|_{L} ; \mu_{I}\left(x^{I}, p_{\bar{I}}\right)-$ плотность меры $\mu \in \Lambda^{n}(L)$ в координатах $\left(x^{I}, p_{\bar{I}}\right): \mu=$ $=\mu_{I} d x^{I}{ }^{\wedge} d p_{\bar{I}}$ с нек-рым выбором аргумента $\arg \mu_{I}$.



\begin{enumerate}

  \setcounter{enumi}{1}

  \item Глобальная теория. Имеет место соотношение

\end{enumerate}



\[

k_{I} \varphi=e^{\frac{i}{h} c_{I J}^{(1)}+i \pi c_{I J}^{(2)}} k_{J} \varphi+O(h), \quad \varphi \in C_{0}^{\infty}\left(U_{I} \cap U_{J}\right)

\]



где



\[

\begin{gathered}

c_{I J}^{(1)}=S_{I}+x^{\bar{I}} p_{\bar{I}}-S_{J}-\left.x^{\bar{J}} p_{J}\right|_{L}, \\

c_{I J}^{(2)}=\frac{1}{2 \pi}\left[\arg \mu_{I}-\arg \mu_{J}-\sum_{k} \arg \lambda_{I J}^{k}-\left|I_{2}\right| \pi\right],

\end{gathered}

\]



причем $\lambda_{I J}^{k}-$ собственные значения матрицы



\begin{center}

\includegraphics[max width=\textwidth]{2023_07_08_69a53b41a3ee340109b6g-1}

\end{center}



I



\[

-\frac{3 \pi}{2}<\arg \lambda_{I J}^{k} \leqslant \frac{\pi}{2}

\]



$c_{I J}^{(1)}, c_{I J}^{(2)}-$ константы. Набор этих констант определяет одномерные классы когомологий $c^{(1)}$ и $c^{(2)}$, к-рые являются препятствиями к существованию глобального канонич. оператора $k$, совпадающего с операторами (1) в картах $U_{I}$. Лагранжево многообразие с мерой, на к-ром $c^{(1)}=0$, $c^{(2)}=0(\bmod 2)$, называется к в а н т о в а н н ы м. На квантованном лагранжевом многообразии определен (с точностью до мультипликативной константы) оператор



\[

k: C_{0}^{\infty}(L) \rightarrow C^{\infty}\left(\mathbb{R}_{x}^{n} \times \mathbb{R}_{+, h}^{1}\right),

\]



к-рый и называется каноническим оператором Маслова.



\begin{enumerate}

  \setcounter{enumi}{2}

  \item Формула коммутации с $1 / h$-псевдодифференциальным оператором. Пусть $\hat{H}=\hat{H}\left(x,-i h \frac{\partial}{\partial x}, h\right)-1 / h$-псевдодифференциальный оператор с главным символом (гамильто- нианом) $H_{0}(x, p)$. Тогда имеет место формула

\end{enumerate}



\[

\hat{H} \circ k(\varphi)=k\left(\left.H_{0}(x, p)\right|_{L} \varphi\right)+O(h) .

\]



Если же потребовать дополнительно, чтобы гамильтониан $H$ обрашался в нуль на лагранжевом многообразии $L$, то формула (4) может быть уточнена:



\[

\hat{H} \circ k(\varphi)=-i h k(\mathscr{P} \varphi)+O\left(h^{2}\right),

\]



где $\mathscr{P}-$ оператор переноса:



\[

\mathscr{P}=\frac{d}{d t}+\frac{1}{2} \frac{L_{d / d t} \mu}{\mu}+H^{\text {sub }},

\]



где $\frac{d}{d t}-$ дифференцирование вдоль гамильтонова векторного поля, отвечающего гамильтониану $H_{0}$, a $H^{\text {sub }}-$ cyбглавный символ оператора $\hat{H}$ :



\[

H^{\mathrm{sub}}(x, p)=i H_{1}(x, p)-\frac{1}{2} \frac{\partial^{2} H_{0}}{\partial x \partial p}(x, p)

\]



причем $H_{1}(x, p)$ - коэффициент при $h$ в разложении полного символа (см. 1/h-псевдодифференциальный оператор).



Операторы $k_{I}$ могут быть модифицированы так, чтобы равенство (2) выполнялось с точностью до $O\left(h^{N}\right)$ при любом $N$. В этом случае может быть получена формула коммутации до любой степени $h$.



\begin{enumerate}

  \setcounter{enumi}{3}

  \item Условия квантования Маслова - условия существования М. к. о. на лагранжевом многообразии $L$ с мерой $\mu$. Они были введены как условия обрашения в нуль двух классов $c^{(1)}, c^{(2)}$ когомологий Чеха канонич. покрытия многообразия $L$ (первое и второе условия квантования). Имеется эквивалентная форма условий квантования в терминах когомологий де $\mathrm{Pa}-$ ма. Классы $c^{(1)}, c^{(2)}$ могут быть заданы следующими замкнутыми дифференциальными формами на $L$ (см. [3]):

\end{enumerate}



\[

\begin{gathered}

c^{(1)}: \omega^{(1)}=p d x, \\

c^{(2)}: \omega^{(2)}=\frac{1}{2 \pi i} d \ln \frac{\mu}{\sigma},

\end{gathered}

\]



где мера $\sigma-$ канонич. мера на лагранжевом многообразии, равная $\sigma=d\left(x^{1}-i p_{1}\right) \wedge \ldots \wedge d\left(x^{n}-i p_{n}\right) ;$ в частности, пара $(L, \sigma)$ всегда удовлетворяет второму условию квантования. Класс $c^{(2)}$ называется и нде кс ным клас с ом Маслов а. Значение класса $c^{(2)}$ на замкнутой кривой называется и нде ксом Масло в а этой кривой. Индекс Маслова кривой $\gamma$ может быть представлен различными способами, как, напр.:



\begin{enumerate}

  \item индекс эллиптич. оператора на кривой

\end{enumerate}



\[

S_{\gamma} \varphi(t)=(\sigma+\mu) \varphi(t)+\frac{\sigma-\mu}{\pi i} \int_{\gamma} \frac{\varphi(\tau) d \tau}{\tau-t}

\]



действующего из пространства $L_{2}(\gamma, \mathbb{C})$ в пространство $L_{2}\left(\gamma, \Lambda^{n}(L)\right)$



\begin{enumerate}

  \setcounter{enumi}{1}

  \item вариация аргумента функции $\mu / \sigma$ вдоль кривой $l$.

\end{enumerate}



Для других типов асимптотических разложений имеются следуюшие конструкции М. к. о.



I. $A-$ а нонический оператор Маслова (см. [1]). Пусть $A: H \rightarrow H-$ неограниченный самосопряженный оператор в (гильбертовом) пространстве $H ; C^{\infty}\left(\mathbb{R}^{n}, H\right)-$ пространство функций; фильтрация в этом пространстве определяется оператором $A$, то есть



\[

E_{k}=\left\{f \in C^{\infty}\left(\mathbb{R}^{n}, H\right) \mid f \in \mathscr{D}\left(A^{k}\right)\right\}

\]



Локальный М. к. о., отвечаюший этой фильтрации:



\[

k_{I} \varphi=F_{p_{\bar{I}} \rightarrow x^{\bar{I}}}^{A}\left\{e^{i S_{I}\left(x^{I}, p_{\bar{I}}\right) A} \sqrt{\mu_{I}\left(x^{I}, p_{\bar{I}}\right)} \varphi\left(x^{I}, p_{\bar{I}}\right)\right\}

\]



\begin{center}

\includegraphics[max width=\textwidth]{2023_07_08_69a53b41a3ee340109b6g-1(1)}

\end{center}



\[

F_{p_{\bar{I}} \rightarrow x^{\bar{I}}}^{A}(f)=\left(\frac{i|A|}{2 \pi}\right)^{n / 2} \int_{\mathbb{R}_{\bar{I}}} e^{i \bar{I}^{\bar{I}} p_{\bar{I}} A} f\left(p_{\bar{I}}\right) d p_{\bar{I}},

\]



где последний интеграл есть интеграл в смысле Бохнера.





\end{document}